% =====================================================================
% Mirante dos Dados — Working Paper #5
% Cirurgia uroginecológica no SUS, 2008–2025: ganhos silenciosos de
% eficiência, desigualdade territorial persistente, choque pandêmico
% e represa cirúrgica.
%
% Autor:  Leonardo Chalhoub (engenharia de dados, análise, redação)
% Versão: 1.0  —  Abril de 2026
% Compilação: pdflatex articles/uropro-saude-publica-2008-2025.tex (2x)
%
% Este Working Paper tem como base o trabalho de especialização em
% Enfermagem de TATIELI (2022). O escopo aqui é amplo (2008–2025),
% foca em dinâmicas longitudinais e incorpora a correção da pipeline
% silver (bug de _ingest_ts identificado e resolvido em abril/2026).
% =====================================================================

\documentclass[12pt,a4paper,oneside]{article}

\usepackage[T1]{fontenc}
\usepackage[utf8]{inputenc}
\usepackage[brazilian]{babel}
\usepackage{lmodern}
\usepackage[a4paper,top=3cm,bottom=2cm,left=3cm,right=2cm]{geometry}
\usepackage{setspace}
\usepackage{indentfirst}
\usepackage{titlesec}
\usepackage{caption}
\usepackage{booktabs}
\usepackage{array}
\usepackage{multirow}
\usepackage{graphicx}
\usepackage[hidelinks]{hyperref}
\usepackage{url}
\usepackage{float}
\usepackage{xcolor}

\onehalfspacing
\setlength{\parindent}{1.25cm}
\setlength{\parskip}{0pt}

\titleformat{\section}{\normalfont\bfseries\large\uppercase}{\thesection}{0.6em}{}
\titleformat{\subsection}{\normalfont\bfseries\normalsize}{\thesubsection}{0.6em}{}
\titlespacing*{\section}{0pt}{18pt}{8pt}
\titlespacing*{\subsection}{0pt}{12pt}{6pt}

\captionsetup[table]{labelfont=bf,name=Tabela,labelsep=endash,
                     position=above,justification=centering,font=small}
\captionsetup[figure]{labelfont=bf,name=Figura,labelsep=endash,
                      position=below,justification=centering,font=small}

\newcommand{\rs}{R\$\,}

\begin{document}

% ---------- CAPA ---------------------------------------------------------
\begin{titlepage}
  \begin{center}
    \vspace*{1.5cm}
    {\small\textsc{Mirante dos Dados \quad\textbar\quad Working Paper n. 5}}\\
    {\small Versão 1.0 \textemdash{} Abril de 2026}

    \vfill

    {\LARGE\bfseries
     CIRURGIA UROGINECOLÓGICA NO SISTEMA
     ÚNICO DE SAÚDE, 2008\textendash 2025:
     GANHOS SILENCIOSOS DE EFICIÊNCIA,
     DESIGUALDADE TERRITORIAL PERSISTENTE,
     CHOQUE PANDÊMICO E REPRESA CIRÚRGICA}

    \vspace{1cm}
    {\large\itshape
     Uro-gynecological surgery under Brazil's Unified Health System,
     2008\textendash 2025: silent efficiency gains, persistent
     territorial inequality, pandemic shock and surgical backlog}

    \vfill

    {\large\textbf{Leonardo Chalhoub}\textsuperscript{1}}\\[6pt]
    {\small\textsuperscript{1}\,Engenheiro de Dados \textemdash{}
     análise quantitativa, pipeline aberto, redação.}\\[6pt]
    {\small\href{mailto:leonardochalhoub@gmail.com}{leonardochalhoub@gmail.com}
     \quad\textbullet\quad
     \href{https://github.com/leonardochalhoub}{github.com/leonardochalhoub}}

    \vfill

    \rule{0.6\linewidth}{0.4pt}\\[6pt]
    {\small\textit{Este artigo é o quinto da série de Working Papers
     do Mirante dos Dados. Tem como base a pesquisa original de
     TATIELI (2022, especialização em Enfermagem). O recorte
     amplia-se para a janela 2008\textendash 2025 (dezessete anos),
     dedicada à dinâmica longitudinal de eficiência, equidade e
     resiliência \textit{intra-vertical}. O cruzamento
     \textit{cross-vertical} com Bolsa Família e Emendas
     Parlamentares é tratado em CHALHOUB (2026d, Working Paper n. 3),
     paper companheiro deste.}}

    \vspace{1cm}
    {\footnotesize
     Disponível em
     \href{https://leonardochalhoub.github.io/mirante-dos-dados-br/}%
     {leonardochalhoub.github.io/mirante-dos-dados-br}}\\[2pt]
    {\footnotesize
     \href{https://github.com/leonardochalhoub/mirante-dos-dados-br}%
     {github.com/leonardochalhoub/mirante-dos-dados-br}}
  \end{center}
\end{titlepage}

% =====================================================================
\section*{Resumo}
\addcontentsline{toc}{section}{Resumo}
% =====================================================================
\begin{singlespace}\small\noindent
\textbf{Objetivo.} Caracterizar a evolução longitudinal do tratamento
cirúrgico da incontinência urinária no Sistema Único de Saúde
brasileiro entre 2008 e 2025, com foco em três dinâmicas: a queda
contínua da permanência hospitalar (ganho silencioso de eficiência);
a desigualdade territorial no acesso, agora visível para todas as
27 unidades da federação após correção do pipeline; e o efeito
disruptivo da pandemia de COVID-19 seguido de uma represa cirúrgica
visível em 2024\textendash 2025.

\textbf{Métodos.} Microdados SIH-AIH-RD do DATASUS para os
procedimentos SIGTAP 0409010499 (via abdominal) e 0409070270
(via vaginal), filtrados a partir do conjunto cru e agregados pela
arquitetura medalhão (\textit{bronze}\,/\,\textit{silver}\,/\,\textit{gold})
no Mirante dos Dados. Despesas deflacionadas pelo IPCA (BCB) para
\rs 2021. Populações estaduais do IBGE/SIDRA. O pipeline foi
auditado e tem reprodutibilidade pública: artefato Delta, código
no GitHub e exportação JSON versionada no repositório.

\textbf{Resultados.} (i)~A permanência hospitalar média na via
vaginal caiu de 2{,}39 dias em 2008 para 1{,}43 dia em 2025
(\textendash 40{,}2\,\%); na via abdominal, de 2{,}47 para
1{,}81 dia (\textendash 26{,}7\,\%). (ii)~A taxa nacional na via
vaginal foi de 5\,956 cirurgias\,/\,ano (média 2015\textendash 2019)
para 2\,539 (média 2020\textendash 2021), uma queda de 57\,\%; em
2024\textendash 2025 voltou para um patamar (8\,966 e 10\,144) acima
do pré-pandêmico, sugerindo represa cirúrgica em fase final de
liquidação. (iii)~A desigualdade territorial é severa: em 2025,
SC e PR alcançam 14{,}71 e 11{,}61 procedimentos por 100\,mil
habitantes; PE, AL e RR ficam abaixo de 0{,}75. (Uma análise
\textit{cross-vertical} dessa desigualdade, cruzando o acesso
cirúrgico com cobertura do Bolsa Família e emendas
parlamentares, é apresentada no paper companheiro CHALHOUB,
2026d, Working Paper n. 3.) (iv)~A mortalidade hospitalar é
estruturalmente baixa (\textless\,0{,}05\,\% \textit{intra-hospital}),
reforçando a tese de que o procedimento é seguro e a barreira é
de oferta, não de risco clínico.

\textbf{Conclusões.} O período examinado (2008\textendash 2025) é
marcado por um \textit{trend} positivo de eficiência hospitalar
(menos dias por internação) que não foi acompanhado por equalização
territorial. A pandemia
expandiu temporariamente a desigualdade no acesso e produziu um
\textit{backlog} cirúrgico cujo escoamento parece estar acontecendo
em 2024\textendash 2025. A integridade do dado público depende não
apenas da publicação dos microdados, mas também da qualidade da
engenharia de dados aplicada \textit{downstream}: documentamos um
defeito específico de pipeline (filtro de \texttt{\_ingest\_ts})
descoberto em abril de 2026, capaz de descartar 73\,\% das linhas
\textit{bronze} \textit{e} 14 das 27 unidades da federação \textit{sem
qualquer aviso}.

\vspace{4pt}
\textbf{Palavras-chave:} SUS; incontinência urinária; cirurgias
eletivas; pandemia COVID-19; permanência hospitalar; desigualdade
territorial; arquitetura medalhão; pipeline de dados; represa
cirúrgica.
\end{singlespace}

% =====================================================================
\section*{Abstract}
\addcontentsline{toc}{section}{Abstract}
% =====================================================================
\begin{singlespace}\small\noindent
\textbf{Objective.} To characterize the longitudinal evolution of
surgical treatment for urinary incontinence in Brazil's Unified
Health System (SUS) between 2008 and 2025, focusing on three
dynamics: the continuous decline in hospital length of stay (a
silent efficiency gain); the territorial inequality of access,
now visible across all 27 federative units after a pipeline
correction; and the disruptive effect of the COVID-19 pandemic
followed by a surgical backlog visible in 2024\textendash 2025.

\textbf{Methods.} SIH-AIH-RD microdata from DATASUS for SIGTAP
procedures 0409010499 (abdominal approach) and 0409070270 (vaginal
approach), filtered from the raw set and aggregated through a
medallion architecture (\textit{bronze}\,/\,\textit{silver}\,/\,\textit{gold})
in the open Mirante dos Dados pipeline. Expenditures deflated by
IPCA (BCB) to \rs 2021. State populations from IBGE/SIDRA.

\textbf{Results.} Vaginal-route length of stay fell from
2.39\,days (2008) to 1.43\,days (2025), a 40.2\,\% reduction.
Volume contracted by 57\,\% in 2020\textendash 2021 vs.
2015\textendash 2019 baseline and rebounded above the pre-pandemic
mean in 2024\textendash 2025, evidencing a backlog draining out.
Per-capita access in 2025 ranges from 14.71/100k (SC) to
0.14/100k (RR), a 100-fold gap. In-hospital mortality is
structurally below 0.05\,\%.

\textbf{Conclusions.} A decade of clinical-efficiency gains was
not matched by territorial equalization. Pandemic shock
temporarily widened access inequality and produced a backlog
visibly draining in 2024\textendash 2025. We document a specific
pipeline defect (\texttt{\_ingest\_ts} filter) discovered in
April 2026 capable of silently dropping 73\,\% of bronze rows
and 14 out of 27 federative units, underlining that public-data
integrity depends on data-engineering quality, not only on
microdata publication.

\vspace{4pt}
\textbf{Keywords:} SUS; urinary incontinence; elective surgery;
COVID-19 pandemic; length of stay; territorial inequality;
medallion architecture; data pipeline; surgical backlog.
\end{singlespace}

% =====================================================================
\section{Introdução}\label{sec:intro}
% =====================================================================

A incontinência urinária (IU) é hoje uma das condições crônicas
mais prevalentes na população adulta brasileira, especialmente
entre mulheres a partir da quarta década de vida. O tratamento
cirúrgico, indicado em casos refratários ao manejo conservador,
é majoritariamente provido pelo Sistema Único de Saúde (SUS) em
hospitais de média e alta complexidade. Esse arranjo institucional
torna o SUS a fonte primária de qualquer caracterização populacional
da prática cirúrgica para IU no país.

A literatura clínica internacional consolidou, ao longo dos últimos
quinze anos, a preferência pelo \textit{sling} sub-uretral
(via vaginal) sobre as suspensões abdominais clássicas (Burch,
Marshall-Marchetti-Krantz). Revisões sistemáticas Cochrane (FORD
\textit{et al.}, 2017) e a declaração de consenso da European
Association of Urology (CHAPPLE \textit{et al.}, 2020) apontam, com
evidência de alto nível, que a abordagem vaginal apresenta menor
tempo cirúrgico, menor permanência hospitalar e taxa de cura
comparável à abdominal. Como mostrou TATIELI (2022), o SUS de
fato consolidou essa preferência: a razão vaginal:abdominal
observada na série SIH-AIH-RD entre 2015 e 2020 estabilizou-se
em torno de 7:1, alinhada à literatura.

O presente artigo \textit{não} é uma replicação. É uma extensão
deliberada em três dimensões. Primeiro, o recorte temporal é
ampliado para 2008\textendash 2025, cobrindo dezessete anos
\textemdash{} quase duas décadas de prática cirúrgica, com o
pós-pandemia já em curso de normalização. Segundo, a análise foca
em três dinâmicas longitudinais que não eram visíveis no recorte
mais curto: a queda silenciosa da permanência hospitalar, a
forma assimétrica da recuperação pós-pandêmica entre as unidades
da federação, e a represa cirúrgica visível em 2024\textendash 2025.
Terceiro, o artigo documenta com transparência um defeito de
\textit{pipeline} descoberto em abril de 2026 e responsável por
ocultar mais de metade das UFs do Brasil em todas as visualizações
prévias da plataforma. Esse último ponto, à primeira vista
metodológico, é o que justifica a existência mesma deste paper:
\textit{os números aqui apresentados refletem a correção da camada
silver; quaisquer divergências em relação às visualizações prévias
da plataforma são causadas por uma correção de engenharia, não por
revisão de fonte.}

A pergunta que organiza o texto é: \textit{ao olharmos quase duas
décadas de prática cirúrgica brasileira para uma única condição
clínica bem caracterizada, o que enxergamos sobre a capacidade
do sistema público de produzir simultaneamente eficiência clínica
e equidade territorial?}

A resposta antecipada é dupla. Por um lado, há ganho real e
contínuo: a permanência hospitalar média caiu mais de 40\,\% na
via vaginal entre 2008 e 2025; a mortalidade \textit{intra-hospitalar}
permanece estruturalmente abaixo de 0{,}05\,\%; a despesa
deflacionada por procedimento manteve-se relativamente estável
em moeda constante. Por outro lado, a desigualdade territorial é
ampla e estável: em 2025, a taxa por 100\,mil habitantes em Santa
Catarina (14{,}71) é mais de cem vezes maior que em Roraima
(0{,}14). A pandemia de COVID-19 não criou essa desigualdade,
mas a explicitou: estados com sistemas hospitalares mais densos
recuperaram-se rapidamente e ainda escoam represa cirúrgica em
2024\textendash 2025; estados periféricos sequer voltaram ao
patamar pré-pandêmico em 2025.

% =====================================================================
\section{Marco Conceitual}\label{sec:teoria}
% =====================================================================

\subsection{Cirurgia eletiva como categoria sensível}

Cirurgias eletivas \textemdash{} aquelas que podem ser
\textit{adiadas sem dano clínico imediato} \textemdash{}
são, por construção, a primeira categoria a ser suprimida quando
um sistema hospitalar é tensionado por uma demanda externa
inelástica. A pandemia de COVID-19 explicitou esse mecanismo de
forma extrema: a realocação compulsória de leitos de enfermaria
e UTI para pacientes com SARS-CoV-2 retirou capacidade do estoque
disponível para procedimentos não-emergenciais (NUNES \textit{et
al.}, 2022). A IU cirúrgica é prototípica dessa categoria: é
quase sempre eletiva, requer hospital de média ou alta
complexidade e demanda anestesia regional ou geral. Quando o
hospital prioriza COVID, a fila de IU não desaparece \textemdash{}
acumula. Esse acúmulo é o que chamaremos de \textit{represa
cirúrgica}.

\subsection{Permanência hospitalar como métrica de eficiência}

A permanência hospitalar média (\textit{average length of stay},
ALOS) é a medida síntese mais simples da intensidade de uso de
recursos por procedimento. Um procedimento que mantém uma boa
taxa de cura mas gradualmente requer menos dias de internação
representa ganho real de produtividade hospitalar: o mesmo leito
pode atender mais pacientes ao longo do ano. No contexto da IU
cirúrgica, esse ganho é compatível com a difusão progressiva de
técnicas minimamente invasivas (\textit{slings} sub-uretrais
TVT/TOT) e com a diminuição da via abdominal pura (que requeria
internação mais longa). A queda da permanência observada nos
microdados é, portanto, a tradução estatística de uma evolução
clínica de fato ocorrida.

\subsection{Equidade territorial em sistema universal}

O princípio constitucional da universalidade do SUS implica
\textit{acesso} universal, não \textit{distribuição} uniforme da
oferta. Estados com maior densidade hospitalar (especialmente Sul
e Sudeste, maiores centros de formação médica) historicamente
concentraram a oferta de procedimentos especializados. Para a IU
cirúrgica, essa concentração não é novidade: TATIELI (2022) já a
documentava no recorte 2015\textendash 2020. O que este Working
Paper acrescenta é que a desigualdade não diminuiu ao longo do
período examinado,
\textit{e} que o pipeline público até abril de 2026 estava
publicando uma versão \textit{ainda mais desigual do que a real},
por força do bug descrito na seção metodológica.

% =====================================================================
\section{Aspectos Metodológicos}\label{sec:metodo}
% =====================================================================

\subsection{Fontes de dados}

A análise utiliza três fontes públicas:
\begin{itemize}
\item \textbf{SIH-AIH-RD} (DATASUS): microdados de Autorizações de
Internação Hospitalar reduzidas, distribuídos via servidor FTP em
arquivos DBC mensais por UF. Período coberto: janeiro/2008 a
dezembro/2025.
\item \textbf{IPCA} (Banco Central do Brasil, série SGS\,433):
Índice de Preços ao Consumidor Amplo para deflação de despesas
para reais constantes (\rs 2021).
\item \textbf{Estimativas populacionais} (IBGE\,/\,SIDRA, tabela
6579): população residente por UF, base anual.
\end{itemize}

\subsection{Pipeline de dados (arquitetura medalhão)}

A arquitetura adotada segue o padrão \textit{lakehouse}
(ARMBRUST \textit{et al.}, 2021), com três camadas:

\begin{itemize}
\item \textbf{Bronze}: ingestão crua via Auto Loader. Os arquivos
DBC originais são mantidos em volumes Unity Catalog; uma camada
\textit{bronze} de Delta Lake armazena cada linha do DBC convertida
para Parquet com tipos exclusivamente \texttt{string}, sem
inferência. Essa decisão (camada \textit{bronze} \textit{string-only})
preserva a fidelidade ao dado de origem e remete decisões de
tipagem para o \textit{silver}.

\item \textbf{Silver}: limpeza, normalização e agregação. Aqui
ocorrem (i) \textit{cast} de tipos para os domínios canônicos
(inteiro, decimal, data, código), (ii) deduplicação e validação,
e (iii) agregação por chave analítica
(uf $\times$ ano $\times$ mês $\times$ \texttt{proc\_rea}
$\times$ caráter $\times$ gestão).

\item \textbf{Gold}: vista analítica colapsada por
(uf $\times$ ano $\times$ \texttt{proc\_rea}), enriquecida com
populações IBGE e deflatores IPCA. É a camada consumida pelo
\textit{front-end} da plataforma e por exportações.
\end{itemize}

\subsection{Sobre o defeito do pipeline silver e sua correção}\label{sub:bug}

Em abril de 2026, durante uma inspeção visual do mapa de
\textit{Distribuição geográfica} da plataforma, observou-se que
diversas unidades da federação apareciam como \textit{sem dado}
para combinações (ano, procedimento) em que o sistema deveria, em
princípio, ter registros. Uma sequência de consultas SQL diretas
ao Delta da camada \textit{silver} revelou que, para o procedimento
de maior volume nacional (0409070270, via vaginal), apenas 13 das
27 UFs apareciam por ano \textemdash{} mesmo nos anos em que o
\textit{bronze} continha registros para todas as 27.

\noindent
A causa raiz, identificada por inspeção do código, foi a seguinte
linha no \textit{notebook silver}:

\begin{singlespace}\small
\begin{verbatim}
latest_ts = bronze.agg(F.max("_ingest_ts")).first()[0]
src = bronze.where(F.col("_ingest_ts") == latest_ts)
\end{verbatim}
\end{singlespace}

\noindent
A intenção declarada (em comentário) era \textit{evitar dupla
contagem se rodaram o ingest mais de uma vez}. O bug está na
hipótese implícita: o filtro pressupõe que \textit{cada} valor
de \texttt{\_ingest\_ts} representa um \textit{snapshot} completo
do \textit{bronze}. Isso não é verdade. O Auto Loader processa
arquivos em micro-\textit{batches} dentro de uma mesma execução,
e \textit{cada} micro-batch recebe seu próprio
\texttt{\_ingest\_ts}. No caso real, os 408 arquivos por UF foram
distribuídos em quatro micro-batches separados por poucos segundos
(00:07:17, 00:07:19, 00:07:21 e 00:07:25 UTC). O filtro
\texttt{== max} preservava \textit{apenas} o último micro-batch,
descartando 73\,\% das linhas \textit{bronze} \textit{e} 14 das 27
UFs \textit{em silêncio}.

\noindent
A correção foi remover o filtro. A idempotência por arquivo é
garantida pelo \textit{checkpoint} do próprio Auto Loader (cada
arquivo é processado uma única vez), o que torna o filtro
\texttt{max(\_ingest\_ts)} não apenas defeituoso mas também
redundante. O \textit{commit} público que aplica a correção é o
\texttt{fa869cf} no repositório do projeto (abril de 2026). Toda a
camada \textit{silver} e \textit{gold} foi regenerada em seguida,
e os números apresentados neste Working Paper já refletem essa
reconstituição. Magnitudes absolutas extraídas da plataforma antes
da correção devem ser consideradas prejudicadas, embora as direções
de tendência se mantenham \textemdash{} o que é um achado
metodológico relevante por si só.

\subsection{Procedimentos analisados}

Foram considerados dois procedimentos da Tabela SIGTAP:
\begin{itemize}
\item \textbf{0409010499} \textemdash{} \textit{Tratamento
Cirúrgico de Incontinência Urinária por Via Abdominal}.
\item \textbf{0409070270} \textemdash{} \textit{Tratamento
Cirúrgico de Incontinência Urinária por Via Vaginal}.
\end{itemize}

O código residual 0409020117 (\textit{Genérico}), embora citado
em estudos prévios, não apresenta registros não-nulos na série
\textit{bronze} obtida pelo \textit{ingest} via FTP do DATASUS, o
que sugere que se trata de código histórico \textit{deprecated}.
Excluímo-lo da análise.

\subsection{Indicadores construídos}

A camada \textit{gold} produz, por \texttt{(uf, ano, proc\_rea)},
os seguintes indicadores que serão usados ao longo do texto:

\begin{itemize}
\item \texttt{n\_aih}: número de internações aprovadas. \textit{Proxy}
direto do volume de cirurgias.
\item \texttt{val\_tot\_2021}: valor total pago em \rs 2021,
deflacionado pelo IPCA.
\item \texttt{val\_per\_aih\_2021}: valor médio pago por AIH em
\rs 2021. Sensível à composição (mix abd/vag) e a reajustes da
SIGTAP.
\item \texttt{dias\_perm\_avg}: permanência hospitalar média
\textit{ponderada por AIH}. Métrica chave deste artigo.
\item \texttt{mortalidade}: taxa de óbito intra-hospitalar
($\#\text{mortes}\,/\,\#\text{AIH}$).
\item \texttt{n\_aih\_por100k}: razão por 100\,mil habitantes,
usando a estimativa IBGE da UF no ano correspondente. Permite
comparação direta entre estados com populações distintas.
\end{itemize}

\subsection{Limitações}

(i) A análise é \textit{descritiva} \textemdash{} não há controle
para confundidores demográficos (estrutura etária da população
feminina, paridade, prevalência subjacente de IU) nem para
heterogeneidade da capacidade hospitalar. (ii) O SIH captura
apenas a parte da prática efetivamente \textit{paga} pelo SUS;
cirurgias particulares e por planos privados não aparecem.
(iii) A definição de IU empregada é a operacional do SIGTAP
(o procedimento foi codificado como tal), não uma definição
clínica externa. (iv) Anos parciais (em geral o ano corrente até
o último mês fechado) são descartados na regra
\textit{full-year-only} aplicada na camada \textit{silver}; isso
mantém a comparabilidade temporal mas exige cautela ao interpretar
o ano final disponível.

% =====================================================================
\section{Resultados}\label{sec:resultados}
% =====================================================================

\subsection{Volumes nacionais e razão entre vias}\label{sub:volumes}

A Tabela~\ref{tab:nat-aih} consolida o volume nacional anual por
via para o período 2008\textendash 2025. A via vaginal predomina
de forma estável: nos anos pré-pandêmicos (2015\textendash 2019),
a razão vaginal:abdominal foi de aproximadamente 7{,}5:1,
ligeiramente maior do que estimativas baseadas em dados parciais
extraídas da plataforma antes da correção do bug descrito na
Seção~\ref{sub:bug}. Em 2025, a razão sobe para 13{,}7:1
\textemdash{} um indício de que a represa cirúrgica em escoamento
está concentrada \textit{quase exclusivamente} na via vaginal,
provavelmente porque é a via majoritária da prática contemporânea.

\begin{table}[H]
\centering
\caption{Volume anual nacional por via cirúrgica (n. de AIH).
Fonte: SIH-AIH-RD (DATASUS), processado pelo Mirante dos Dados.}
\label{tab:nat-aih}
\small
\begin{tabular}{rrrr}
\toprule
\textbf{Ano} & \textbf{Abdominal} & \textbf{Vaginal} & \textbf{Razão V:A} \\
\midrule
2008 & 2\,074 & 6\,558  &  3{,}2:1 \\
2009 & 1\,755 & 5\,888  &  3{,}4:1 \\
2010 & 1\,542 & 6\,324  &  4{,}1:1 \\
2011 & 1\,388 & 6\,085  &  4{,}4:1 \\
2012 & 1\,216 & 5\,982  &  4{,}9:1 \\
2013 & 1\,046 & 5\,494  &  5{,}3:1 \\
2014 & 1\,032 & 5\,785  &  5{,}6:1 \\
2015 & 1\,020 & 5\,722  &  5{,}6:1 \\
2016 &   789  & 5\,467  &  6{,}9:1 \\
2017 &   760  & 5\,419  &  7{,}1:1 \\
2018 &   760  & 5\,555  &  7{,}3:1 \\
2019 &   759  & 5\,959  &  7{,}9:1 \\
2020 &   384  & 2\,566  &  6{,}7:1 \\
2021 &   323  & 2\,511  &  7{,}8:1 \\
2022 &   646  & 5\,578  &  8{,}6:1 \\
2023 &   821  & 6\,785  &  8{,}3:1 \\
2024 &   863  & 8\,966  & 10{,}4:1 \\
2025 &   740  & 10\,144 & 13{,}7:1 \\
\bottomrule
\end{tabular}
\end{table}

Três fenômenos saltam à vista. Primeiro, a via abdominal vinha
caindo gradualmente \textit{antes} da pandemia: de 2\,074 em 2008
para 759 em 2019, uma redução de 63\,\% em 11 anos. Esse declínio
é consistente com a difusão internacional dos \textit{slings} e
com a substituição clínica da via abdominal pela vaginal. Segundo,
ambas as vias colapsaram em 2020\textendash 2021, mas a magnitude
relativa foi diferente: a abdominal caiu mais (\textendash 49{,}4\,\%
em 2020) que a vaginal (\textendash 56{,}9\,\%); contudo, em base
absoluta, a vaginal perdeu 3\,393 procedimentos\,/\,ano enquanto
a abdominal perdeu 375. Terceiro, a recuperação pós-pandêmica é
\textit{heterogênea entre vias}: a vaginal não apenas voltou ao
patamar pré-pandêmico em 2022\textendash 2023 (5\,578 e 6\,785,
contra 5\,959 em 2019) como o ultrapassou em 2024\textendash 2025
(8\,966 e 10\,144), enquanto a abdominal estabilizou-se em torno
de 800\,/\,ano \textemdash{} um patamar inferior ao último ano
pré-pandêmico (759), mas dentro da banda de variação histórica
observada desde 2017.

A Figura~\ref{fig:evolucao-vol} ilustra o padrão. O salto de
2024\textendash 2025 na via vaginal é a evidência mais direta
disponível de que a represa cirúrgica acumulada durante a pandemia
está em fase final de escoamento.

\begin{figure}[H]
\centering
\fbox{\parbox[c][6cm][c]{0.85\linewidth}{\centering
\textit{[Figura disponível na versão interativa da plataforma.}\\
\textit{Reprodução estática deste Working Paper omitida nesta
versão; ver gráfico empilhado por via cirúrgica em
leonardochalhoub.github.io/mirante-dos-dados-br/incontinencia-urinaria.]}}}
\caption{Evolução nacional do número de cirurgias por via,
2008\textendash 2025. Barras empilhadas por via (abdominal, vaginal).
A queda 2020\textendash 2021 e o pico 2024\textendash 2025 na via
vaginal são as duas anomalias mais evidentes.}
\label{fig:evolucao-vol}
\end{figure}

\subsection{Permanência hospitalar: a queda silenciosa}\label{sub:permanencia}

Este é, de longe, o achado mais relevante deste Working Paper para
o debate sobre \textit{eficiência} do gasto público em saúde. A
permanência hospitalar média \textit{caiu continuamente em ambas
as vias} ao longo de todo o período, sem interrupção pela pandemia
(que de fato \textit{acelerou} a queda em 2020\textendash 2021).
A Tabela~\ref{tab:nat-perm} sintetiza.

\begin{table}[H]
\centering
\caption{Permanência hospitalar média nacional, em dias por AIH.
Ponderação: número de AIH em cada ano. Fonte: SIH-AIH-RD.}
\label{tab:nat-perm}
\small
\begin{tabular}{rrr}
\toprule
\textbf{Ano} & \textbf{Abdominal} & \textbf{Vaginal} \\
\midrule
2008 & 2{,}47 & 2{,}39 \\
2009 & 2{,}44 & 2{,}41 \\
2010 & 2{,}47 & 2{,}37 \\
2011 & 2{,}13 & 2{,}29 \\
2012 & 2{,}22 & 2{,}22 \\
2013 & 2{,}22 & 2{,}02 \\
2014 & 2{,}48 & 1{,}97 \\
2015 & 2{,}24 & 1{,}92 \\
2016 & 2{,}16 & 1{,}90 \\
2017 & 2{,}15 & 1{,}83 \\
2018 & 2{,}23 & 1{,}77 \\
2019 & 2{,}14 & 1{,}86 \\
2020 & 1{,}96 & 1{,}68 \\
2021 & 1{,}94 & 1{,}55 \\
2022 & 1{,}86 & 1{,}52 \\
2023 & 1{,}79 & 1{,}53 \\
2024 & 1{,}86 & 1{,}45 \\
2025 & 1{,}81 & 1{,}43 \\
\midrule
\textbf{$\Delta$ 2008$\rightarrow$2025} & \textbf{\textendash 26{,}7\,\%} & \textbf{\textendash 40{,}2\,\%} \\
\bottomrule
\end{tabular}
\end{table}

Em 2008, a permanência média na via vaginal era de 2{,}39 dias.
Em 2025, é de 1{,}43 \textemdash{} uma redução de 40{,}2\,\% em
17 anos, ou pouco mais de 1\,\% ao ano em média geométrica. Na
via abdominal, a queda foi de 2{,}47 para 1{,}81 dia
(\textendash 26{,}7\,\%), também substantiva. Esses números têm
três implicações imediatas.

\textbf{Primeira, e mais óbvia: o sistema produz hoje, com o
mesmo leito-dia, mais cirurgias do que em 2008.} Em ordem de
grandeza: 1\,000 leitos-dia em 2008 entregavam aproximadamente
418 procedimentos vaginais (1\,000 / 2{,}39); em 2025 entregam
699 (1\,000 / 1{,}43) \textemdash{} um ganho de produtividade
hospitalar de 67\,\%. Isso é um \textit{economia real de recurso
público} mantida silenciosamente fora do radar do debate
político-orçamentário.

\textbf{Segunda: a queda é compatível, em ritmo e direção, com
a difusão internacional dos \textit{slings} sub-uretrais
minimamente invasivos.} A literatura clínica (FORD \textit{et
al.}, 2017; CHAPPLE \textit{et al.}, 2020) reporta permanência
de 1\textendash 2 dias para TVT/TOT em centros bem treinados.
O Brasil, em 2025, está dentro dessa banda \textemdash{}
embora, como veremos, com forte heterogeneidade entre UFs.

\textbf{Terceira: a pandemia não estancou nem reverteu o
\textit{trend}.} Os anos 2020\textendash 2021 mostram, ao
contrário, \textit{aceleração} da queda (de 1{,}86 em 2019 para
1{,}68 em 2020 e 1{,}55 em 2021 na via vaginal). A interpretação
plausível: durante a pandemia, apenas casos com indicação mais
clara, equipes mais treinadas e logística hospitalar mais
azeitada mantiveram-se em operação \textemdash{} um efeito de
\textit{seleção de qualidade} no \textit{numerador} do volume,
que reduz a permanência média mesmo sem mudança técnica adicional.
A relevância desse efeito persiste no pós-pandemia (1{,}52 em
2022, 1{,}43 em 2025): a pandemia parece ter atuado como um
catalisador irreversível de eficiência.

\subsection{Despesa pública e custo médio por AIH}\label{sub:despesa}

A despesa total deflacionada (Tabela~\ref{tab:nat-val}) acompanha
fielmente o volume, com uma exceção interessante: o salto entre
2012 e 2013, em que o valor médio por AIH (vaginal) sobe de
\rs 386 para \rs 677 sem queda equivalente no volume. Esse salto
não é movimento real de mercado; é o reajuste da Tabela SIGTAP
do procedimento (BRASIL, 2024). De 2013 em diante, o valor por
AIH em moeda constante \textit{deflacionada} cai gradualmente
\textemdash{} ou seja, a tabela do SUS perdeu poder real ao
longo dos anos seguintes, sendo apenas parcialmente reajustada
nominal.

\begin{table}[H]
\centering
\caption{Despesa pública nacional total e por AIH, em \rs 2021,
deflacionada pelo IPCA. Procedimento vaginal (0409070270).}
\label{tab:nat-val}
\small
\begin{tabular}{rrrr}
\toprule
\textbf{Ano} & \textbf{Total (R\$\,2021)} & \textbf{Por AIH (R\$\,2021)} & \textbf{n.\,AIH} \\
\midrule
2008 & 2\,181\,785 & 332{,}69 &  6\,558 \\
2010 & 2\,404\,941 & 380{,}29 &  6\,324 \\
2012 & 2\,309\,820 & 386{,}13 &  5\,982 \\
2013 & 3\,719\,105 & 676{,}94 &  5\,494 \\
2015 & 3\,523\,114 & 615{,}71 &  5\,722 \\
2017 & 2\,773\,664 & 511{,}84 &  5\,419 \\
2019 & 2\,951\,955 & 495{,}38 &  5\,959 \\
2020 & 1\,093\,490 & 426{,}15 &  2\,566 \\
2021 &   966\,605 & 384{,}95 &  2\,511 \\
2022 & 2\,071\,890 & 371{,}44 &  5\,578 \\
2023 & 2\,631\,654 & 387{,}86 &  6\,785 \\
2024 & 4\,995\,926 & 557{,}21 &  8\,966 \\
2025 & 4\,328\,034 & 426{,}66 & 10\,144 \\
\bottomrule
\end{tabular}
\end{table}

Em 2025, o gasto público nacional na via vaginal foi de
\rs 4{,}33\,milhões (em \rs 2021), o maior valor da série
em volume real. Por AIH, contudo, o valor foi de \rs 426{,}66
\textemdash{} 16\,\% \textit{abaixo} do pico de 2015
(\rs 615{,}71) e 14\,\% \textit{abaixo} de 2019 (\rs 495{,}38).
O sistema fez \textit{mais cirurgias} pagando \textit{menos por
cirurgia} em moeda constante. Combinada com a redução de
permanência da seção anterior, a leitura é:
\textbf{o gasto cresceu em volume mas o custo unitário caiu}.

\subsection{Geografia: concentração e desigualdade}\label{sub:geografia}

A Tabela~\ref{tab:uf-100k} apresenta a taxa por 100\,mil habitantes
em 2025 para a via vaginal, ordenada do maior ao menor.

\begin{table}[H]
\centering
\caption{Acesso à cirurgia vaginal de IU em 2025: AIH por
100\,mil habitantes, por UF. População IBGE/SIDRA tabela 6579.}
\label{tab:uf-100k}
\small
\begin{tabular}{lrrr}
\toprule
\textbf{UF} & \textbf{AIH 2025} & \textbf{Pop. (M)} & \textbf{AIH/100k} \\
\midrule
SC          & 1\,204 &  8{,}19 & 14{,}71 \\
PR          & 1\,381 & 11{,}89 & 11{,}61 \\
GO          &   621  &  7{,}42 &  8{,}37 \\
MS          &   198  &  2{,}92 &  6{,}77 \\
SP          & 2\,847 & 46{,}08 &  6{,}18 \\
ES          &   237  &  4{,}13 &  5{,}74 \\
RS          &   638  & 11{,}23 &  5{,}68 \\
TO          &    87  &  1{,}59 &  5{,}48 \\
MT          &   160  &  3{,}89 &  4{,}11 \\
MG          &   816  & 21{,}39 &  3{,}81 \\
AM          &   155  &  4{,}32 &  3{,}59 \\
MA          &   251  &  7{,}02 &  3{,}58 \\
RO          &    59  &  1{,}75 &  3{,}37 \\
PI          &   111  &  3{,}38 &  3{,}28 \\
BA          &   457  & 14{,}87 &  3{,}07 \\
CE          &   195  &  9{,}27 &  2{,}10 \\
DF          &    62  &  3{,}00 &  2{,}07 \\
RN          &    69  &  3{,}46 &  2{,}00 \\
AC          &    17  &  0{,}88 &  1{,}92 \\
RJ          &   317  & 17{,}22 &  1{,}84 \\
SE          &    39  &  2{,}30 &  1{,}70 \\
PB          &    67  &  4{,}16 &  1{,}61 \\
AP          &    12  &  0{,}81 &  1{,}49 \\
PE          &    70  &  9{,}56 &  0{,}73 \\
PA          &    63  &  8{,}71 &  0{,}72 \\
AL          &    10  &  3{,}22 &  0{,}31 \\
RR          &     1  &  0{,}74 &  0{,}14 \\
\bottomrule
\end{tabular}
\end{table}

Os números são notáveis. Santa Catarina, no topo, atinge 14{,}71
procedimentos por 100\,mil habitantes em 2025. Roraima, na ponta
inferior, registra 0{,}14 \textemdash{} uma diferença de cerca
de 100 vezes. Pernambuco, com quase 10 milhões de habitantes,
realiza apenas 70 cirurgias vaginais no ano, o que corresponde a
0{,}73\,/\,100k, taxa similar à de Roraima e oito vezes inferior
à de SP.

Algumas observações qualitativas:

\textbf{(i) A liderança do Sul é estrutural, não casual.} SC e
PR ocupam topo do \textit{ranking} sistematicamente desde 2015,
e a vantagem só se ampliou nos anos pós-pandêmicos: a comparação
entre as médias 2015\textendash 2019 e 2022\textendash 2025
mostra crescimento do volume médio em SC ($\Delta$ +138\,\%)
e PR ($\Delta$ +84\,\%) acima do crescimento médio nacional.

\textbf{(ii) A região Sudeste é heterogênea.} SP é forte
(6{,}18/100k), MG modesta (3{,}81), RJ baixa (1{,}84). RJ chama
particular atenção porque, com 17 milhões de habitantes, realiza
em 2025 apenas 317 cirurgias vaginais \textemdash{} aproximadamente
um terço do volume \textit{absoluto} de SC (que tem menos da
metade da população).

\textbf{(iii) A região Nordeste apresenta um vácuo de oferta
mesmo em estados populosos.} BA (3{,}07/100k), PE (0{,}73),
CE (2{,}10), AL (0{,}31) operam taxas muito inferiores às
do Sul. A diferença não é explicada por demanda demográfica
(estrutura etária da população feminina nordestina é, se
algo, mais favorável à incidência de IU que a do Sul), mas
provavelmente por densidade hospitalar especializada e
treinamento médico em cirurgia uroginecológica.

\textbf{(iv) Estados periféricos têm volumes irrelevantes em
qualquer ano.} RR, AP, AL e SE operam taxas abaixo de 2/100k
em toda a série \textemdash{} o que, em termos absolutos,
significa entre uma e algumas dezenas de cirurgias por ano.
Para essas UFs, a hipótese mais econômica é que o tratamento
cirúrgico é deslocado para grandes centros de referência fora
do estado, via Tratamento Fora do Domicílio (TFD) do SUS.

\subsection{O choque pandêmico e a forma da recuperação}\label{sub:covid}

A Tabela~\ref{tab:uf-pre-post} compara, para a via vaginal, as
médias anuais por UF nos três regimes: pré-pandemia
(2015\textendash 2019), pandemia (2020\textendash 2021) e
pós-pandemia (2022\textendash 2025). A coluna $\Delta$pré$\rightarrow$pós
captura a variação relativa entre o regime pré e o pós, e é a
quantidade-chave para detectar represa cirúrgica.

\begin{table}[H]
\centering
\caption{Heterogeneidade da resposta pandêmica por UF: média anual
de cirurgias vaginais (n. AIH) em três regimes.
$\Delta$pré$\rightarrow$pós: variação relativa entre 2015\textendash 19
e 2022\textendash 25.}
\label{tab:uf-pre-post}
\small
\begin{tabular}{lrrrr}
\toprule
\textbf{UF} & \textbf{2015\textendash 19} & \textbf{2020\textendash 21} & \textbf{2022\textendash 25} & \textbf{$\Delta$pré$\rightarrow$pós} \\
\midrule
SP & 1\,964 & 1\,032 & 2\,410 & +22{,}7\,\% \\
PR &   522  &   184  &   958  & +83{,}5\,\% \\
RS &   448  &   199  &   530  & +18{,}3\,\% \\
MG &   440  &   145  &   647  & +47{,}0\,\% \\
GO &   382  &   102  &   437  & +14{,}4\,\% \\
SC &   290  &   181  &   691  & +138{,}3\,\% \\
RJ &   239  &    79  &   231  & \textendash 3{,}3\,\% \\
BA &   202  &    84  &   320  & +58{,}4\,\% \\
MS &   137  &    28  &   197  & +43{,}8\,\% \\
MA &   120  &    76  &   290  & +141{,}7\,\% \\
CE &   117  &   105  &   194  & +65{,}8\,\% \\
PE &   108  &    33  &    70  & \textendash 35{,}2\,\% \\
DF &   101  &    56  &    78  & \textendash 22{,}8\,\% \\
MT &    90  &    45  &   156  & +73{,}3\,\% \\
ES &    87  &    36  &   119  & +36{,}8\,\% \\
PA &    66  &    28  &    64  & \textendash 3{,}0\,\% \\
AM &    55  &    19  &   113  & +105{,}5\,\% \\
AP &    48  &    13  &     9  & \textendash 81{,}2\,\% \\
PE/AL/RR \textit{et al.} & \multicolumn{4}{c}{\textit{(ver discussão)}} \\
\bottomrule
\end{tabular}
\end{table}

A figura emergente é a de \textit{três tipos} de UF.

\textbf{Tipo A \textemdash{} represa em escoamento.} Estados
onde o volume pós-pandêmico ultrapassou \textit{substancialmente}
o pré-pandêmico, indicando que o sistema absorveu a fila
acumulada e segue produzindo acima da tendência prévia. Exemplos:
SC (+138\,\%), MA (+142\,\%), AM (+106\,\%), PR (+84\,\%),
MT (+73\,\%), CE (+66\,\%), BA (+58\,\%), MG (+47\,\%). Em
SC e PR, especificamente, o salto é de tal magnitude que sugere
uma combinação de represa \textit{e} expansão estrutural da
oferta (mais hospitais, mais equipes treinadas).

\textbf{Tipo B \textemdash{} retorno à tendência.} Estados em
que o volume pós-pandêmico está dentro de uma banda de
$\pm$\,25\,\% do pré, sem evidência clara de represa nem de
expansão. Exemplos: SP (+23\,\%), RS (+18\,\%), GO (+14\,\%),
ES (+37\,\%), MS (+44\,\%). Esses estados parecem ter
absorvido o choque e voltado a um regime parecido com o anterior.

\textbf{Tipo C \textemdash{} recuperação incompleta ou regressão.}
Estados em que o volume pós-pandêmico está \textit{abaixo} do
pré-pandêmico, sugerindo perda \textit{permanente} de
capacidade ou de fluxo de pacientes para tratamento cirúrgico
de IU. Exemplos: RJ (\textendash 3\,\%), PA (\textendash 3\,\%),
DF (\textendash 23\,\%), PE (\textendash 35\,\%), AP
(\textendash 81\,\%). RJ e PE são particularmente preocupantes
porque são estados populosos e com tradição médica
\textemdash{} a redução não pode ser explicada por escassez
populacional. Hipóteses qualitativas: (a)~degradação do parque
hospitalar dedicado; (b)~migração de pacientes para particular;
(c)~reorganização da rede que deixou IU sem prioridade.

A magnitude bruta do choque pandêmico (a queda de
2020\textendash 2021) é, ao contrário do padrão de recuperação,
relativamente \textit{homogênea} entre UFs. A maioria dos
estados teve queda entre 50 e 75\,\% em relação ao baseline
2015\textendash 2019. A heterogeneidade emerge \textit{na
recuperação}, não no choque. Esse é um achado relevante para
política pública: o sistema federal foi tensionado de forma
parecida em todos os lugares, mas a capacidade local de
\textit{escoamento da represa} variou enormemente.

\subsection{Mortalidade hospitalar: piso de segurança}\label{sub:mort}

Por fim, um dado positivo e confortavelmente estável: a
mortalidade hospitalar associada a esses procedimentos é
estruturalmente baixíssima. Considerando a série completa
2008\textendash 2025, registramos no total apenas alguns casos
isolados de óbito intra-hospitalar entre as dezenas de milhares
de internações. A taxa nacional anual oscila entre 0\,\% e
0{,}05\,\%; a maior parte dos anos a mortalidade observada na
camada \textit{gold} é literalmente zero. Isso é compatível com
o que a literatura clínica reporta sobre o procedimento: trata-se
de cirurgia de baixo risco quando bem indicada e bem conduzida.

\textbf{Implicação principal:} as barreiras de acesso documentadas
na Seção~\ref{sub:geografia} \textit{não são barreiras de risco
clínico}. Não há justificativa epidemiológica plausível para a
diferença de 100 vezes entre SC e RR no acesso à cirurgia. A
barreira é de \textit{oferta}: há ou não há equipe treinada, há
ou não há hospital com leito disponível, há ou não há fluxo de
referência funcionando. Quando há, a operação é segura e curta.

% =====================================================================
\section{Discussão}\label{sec:discussao}
% =====================================================================

\subsection{Eficiência clínica vs. acesso desigual: a tensão central}

A leitura conjunta dos resultados expõe uma tensão característica
do SUS. De um lado, a Seção~\ref{sub:permanencia} mostra um sistema
\textit{tecnologicamente convergente}: a permanência hospitalar
em 2025 está dentro da banda de centros internacionais de excelência
para o procedimento. Onde a cirurgia é realizada, ela é
\textit{tão eficiente quanto a melhor literatura comparável}. De
outro lado, a Seção~\ref{sub:geografia} mostra um sistema
\textit{territorialmente fragmentado}: a oferta cirúrgica está
extremamente concentrada em alguns estados, deixando regiões
inteiras com taxas de acesso de menos de um décimo do líder.

A combinação dessas duas leituras sugere que o problema da IU
cirúrgica no Brasil em 2025 \textit{não é tecnológico}, é
\textit{logístico}. As equipes brasileiras já fazem o procedimento
bem; o desafio é \textit{onde} ele é feito. A política pública
relevante seria, portanto, mais voltada à expansão da rede
hospitalar capacitada (e da oferta de Tratamento Fora do
Domicílio onde isso for inviável a curto prazo) do que a
investimentos em pesquisa ou inovação técnica.

\subsection{Lições da pandemia: o que aprendemos sobre fragilidade}

A pandemia teve, sob o aspecto narrativo deste artigo, duas
funções distintas. Como \textit{choque}, demonstrou que a
oferta cirúrgica eletiva é altamente sensível a competição por
recursos hospitalares: bastou cerca de um ano de tensionamento
para reduzir o volume nacional pela metade. Como
\textit{revelador}, mostrou que alguns sistemas estaduais
(notadamente SC e PR) eram capazes de responder à represa
acumulada com expansões significativas, enquanto outros (RJ,
PE, AP) não conseguiram sequer voltar ao patamar anterior. A
pandemia funcionou, em outras palavras, como um \textit{teste
de carga} que quantificou diferenças de robustez já
estruturalmente presentes.

A literatura internacional sobre cirurgias eletivas durante e
após a COVID-19 (NUNES \textit{et al.}, 2022) reportou padrões
parecidos em outros países: choque homogêneo, recuperação
heterogênea. O caso brasileiro adiciona uma observação
específica: a heterogeneidade da recuperação no Brasil não é
explicada por uma única dimensão (riqueza estadual, densidade
médica ou tamanho populacional) tomada isoladamente. SC tem PIB
per capita semelhante ao de SP, mas teve crescimento muito
maior em volume cirúrgico pós-pandemia; RJ, com PIB e densidade
médica historicamente altos, regrediu. A explicação parece exigir
fatores institucionais locais (organização da rede de média
complexidade, programas estaduais de capacitação, fluxo de
referência) que não estão capturados na fonte SIH.

\subsection{Sobre o defeito do pipeline e a credibilidade do dado}

Vale a pena tratar explicitamente um ponto incômodo: as tabelas
deste Working Paper \textit{não coincidem} com tabelas e
visualizações prévias da plataforma Mirante. A diferença, contudo,
não é pequena: na tabela de volumes por UF, por exemplo, mais de
metade dos estados aparecia com volume zero em vários anos nas
versões pré-correção \textemdash{} um padrão que era, afinal,
artefato do filtro \texttt{\_ingest\_ts} descrito na
Seção~\ref{sub:bug}, não realidade dos dados.

Isso obriga uma reflexão honesta. \textit{Dados públicos de
qualidade} não são garantidos pela publicação do dado bruto.
Microdados SIH-AIH-RD foram publicados pelo DATASUS por todos
os anos da série, completos, com cobertura nacional plena. O
problema esteve \textit{a jusante}: o pipeline analítico
introduziu uma decisão algorítmica defeituosa que filtrava
silenciosamente dois terços do material bruto. Visualizações
subsequentes pareciam corretas porque as omissões eram
\textit{consistentes} dentro do que era exibido; nenhum alerta
automático de qualidade detectou a perda.

A lição prática, e que pretende ser \textit{de fato útil} a
quem opera pipelines analíticos sobre dados públicos brasileiros:

\begin{enumerate}
\item Evitar filtros do tipo \texttt{== max(\_ingest\_ts)} sobre
fontes ingeridas por Auto Loader. A idempotência por arquivo já
é garantida pelo \textit{checkpoint}; o filtro adicional não é
defensivo, é destrutivo.
\item Implementar \textit{contagens-sentinela}: para fontes onde
sabe-se \textit{a priori} a cardinalidade esperada (como o
``27 UFs'' brasileiras), validar essa cardinalidade explicitamente
em cada camada e quebrar o \textit{job} se for menor.
\item Manter a camada \textit{bronze} como \textit{string-only}:
isso preserva o dado de origem na sua forma mais auditável e
facilita a reconstituição quando algum bug \textit{a jusante} é
descoberto. Foi exatamente o que permitiu a correção rápida deste
caso.
\end{enumerate}

A versão atual da plataforma Mirante dos Dados implementa as
três recomendações acima. O \textit{commit} de correção é
\texttt{fa869cf} no repositório aberto; a regeneração das camadas
\textit{silver} e \textit{gold} foi orquestrada via Databricks
Asset Bundle. O dado consumido pelo \textit{front-end} a partir
da publicação deste Working Paper representa o estado corrigido.

\subsection{Limitações e agenda de pesquisa}

Este artigo é descritivo. Três extensões são naturais.
\textbf{Primeira:} estratificação por sexo e faixa etária, que
exigiria descer para a granularidade de paciente (campo
\texttt{IDADE} e \texttt{SEXO} já presentes no SIH-RD, ainda
não incorporados ao \textit{gold}).
\textbf{Segunda:} cruzamento com a base CNES (estabelecimentos)
para tipologizar a oferta hospitalar (porte, vínculo, sub-rede)
e tornar tratável a hipótese de que a heterogeneidade de
recuperação é explicada pela tipologia da rede local.
\textbf{Terceira:} análise de sobrevida e reoperação a partir
do enlace longitudinal de AIHs por paciente, possível via
campo \texttt{INSC\_PN} mas dependente de tratamento adequado de
identificadores.

% =====================================================================
\section{Considerações Finais}\label{sec:conclusao}
% =====================================================================

A cirurgia uroginecológica no SUS, examinada na escala de
dezessete anos contínuos (2008\textendash 2025), conta uma
história mista.
Ganhamos eficiência clínica continuamente: a permanência caiu
40\,\% na via vaginal e 27\,\% na abdominal, sem inflexão
adversa. Mantivemos segurança: a mortalidade hospitalar é
desprezível em todo período. Conseguimos absorver, mesmo que de
forma desigual, o choque mais agudo do sistema sanitário
brasileiro recente \textemdash{} a pandemia COVID-19 \textemdash{}
e estamos, em 2024\textendash 2025, escoando o estoque de fila
acumulado.

Falhamos, simultaneamente, em um aspecto estruturalmente sensível:
a desigualdade territorial de acesso. A diferença de 100 vezes
entre Santa Catarina e Roraima no acesso por 100\,mil habitantes
não tem qualquer correspondência epidemiológica. É puramente
uma desigualdade de oferta. A pandemia tornou essa desigualdade
ainda mais visível ao mostrar que alguns estados são capazes de
\textit{ampliar a oferta sob estresse} enquanto outros não.

Em termos do dado público em si, este artigo registra um achado
metodológico de relevância \textit{a jusante} dos microdados:
defeitos sutis de \textit{pipeline} podem ocultar mais de
metade da realidade nacional sem qualquer aviso. O acesso público
ao dado bruto é necessário, mas não suficiente, para
caracterizações populacionais robustas. A engenharia de dados
analítica é parte integral da infraestrutura de dado público,
e merece os mesmos padrões de auditoria que aplicamos a outros
componentes da política pública.

% =====================================================================
\section*{Referências}\label{sec:ref}
\addcontentsline{toc}{section}{Referências}
% =====================================================================
\begin{singlespace}\small
\noindent\hangindent=1.25cm\hangafter=1
ARMBRUST, M.; GHODSI, A.; XIN, R.; ZAHARIA, M. Lakehouse: a new
generation of open platforms that unify data warehousing and
advanced analytics. In: \textbf{11th CIDR}, 2021.\par\smallskip

\noindent\hangindent=1.25cm\hangafter=1
BANCO CENTRAL DO BRASIL (BCB). \textbf{Sistema Gerenciador de
Séries Temporais}, série 433: IPCA. Disponível em:
\url{https://api.bcb.gov.br/dados/serie/bcdata.sgs.433/dados}.
Acesso em: abr. 2026.\par\smallskip

\noindent\hangindent=1.25cm\hangafter=1
BRASIL. Ministério da Saúde. \textbf{Sistema de Informações
Hospitalares do SUS (SIH-SUS)}. Microdados. Brasília: DATASUS.
Disponível em:
\url{https://datasus.saude.gov.br/transferencia-de-arquivos/}.
Acesso em: abr. 2026.\par\smallskip

\noindent\hangindent=1.25cm\hangafter=1
BRASIL. Ministério da Saúde. \textbf{Tabela de Procedimentos,
Medicamentos e OPM do SUS \textemdash{} SIGTAP}. Brasília:
DATASUS, 2024.\par\smallskip

\noindent\hangindent=1.25cm\hangafter=1
CHALHOUB, L. Emendas Parlamentares no Orçamento Federal
Brasileiro (2015\textendash 2025): distribuição espacial,
execução orçamentária e efeitos das mudanças institucionais
recentes. \textbf{Mirante dos Dados \textemdash{} Working Paper
n. 1}, abr. 2026a.\par\smallskip

\noindent\hangindent=1.25cm\hangafter=1
CHALHOUB, L. Programa Bolsa Família, Auxílio Brasil e Novo
Bolsa Família (2013\textendash 2025): transformações
institucionais, expansão da cobertura e desigualdade territorial.
\textbf{Mirante dos Dados \textemdash{} Working Paper n. 2},
abr. 2026b.\par\smallskip

\noindent\hangindent=1.25cm\hangafter=1
CHAPPLE, C. R.; CRUZ, F.; DESCHAMPS, C.; HAYLEN, B. T.
\textit{et al.} Consensus statement of the European Association
of Urology on the surgical treatment of urinary incontinence in
women. \textbf{European Urology}, v. 78, n. 5,
p. 643\textendash 656, 2020.\par\smallskip

\noindent\hangindent=1.25cm\hangafter=1
CHALHOUB, L. Acesso desigual: cirurgia uroginecológica no SUS
como indicador de pobreza estrutural e os limites da
compensação fiscal por emendas parlamentares
(2008\textendash 2025). \textbf{Mirante dos Dados \textemdash{}
Working Paper n. 3}, abr. 2026d.\par\smallskip

\noindent\hangindent=1.25cm\hangafter=1
FORD, A. A.; ROGERSON, L.; CODY, J. D.; OGAH, J. Mid-urethral
sling operations for stress urinary incontinence in women.
\textbf{Cochrane Database of Systematic Reviews}, n. 7,
CD006375, 2017.\par\smallskip

\noindent\hangindent=1.25cm\hangafter=1
HAYLEN, B. T.; DE RIDDER, D.; FREEMAN, R. M. \textit{et al.}
An International Urogynecological Association
(IUGA)\,/\,International Continence Society (ICS) joint report
on the terminology for female pelvic floor dysfunction.
\textbf{Neurourology and Urodynamics}, v. 29, n. 1,
p. 4\textendash 20, 2010.\par\smallskip

\noindent\hangindent=1.25cm\hangafter=1
INSTITUTO BRASILEIRO DE GEOGRAFIA E ESTATÍSTICA (IBGE).
\textbf{SIDRA \textemdash{} Sistema IBGE de Recuperação Automática}.
Tabela 6579: Estimativas anuais da população residente.
Disponível em: \url{https://sidra.ibge.gov.br/tabela/6579}.
Acesso em: abr. 2026.\par\smallskip

\noindent\hangindent=1.25cm\hangafter=1
NUNES, B. P.; FLORES, T. R.; MIELKE, G. I. \textit{et al.} The
COVID-19 pandemic and elective surgeries in the Brazilian
Unified Health System (SUS). \textbf{Public Health}, v. 211,
p. 109\textendash 115, 2022.\par\smallskip

\noindent\hangindent=1.25cm\hangafter=1
TATIELI DA SILVA. \textbf{Análise Descritiva da Cirurgia para
Tratamento da Incontinência Urinária no SUS, 2015\textendash 2020}.
Trabalho de Conclusão de Curso (Especialização em Enfermagem),
2022.\par\smallskip

\noindent\hangindent=1.25cm\hangafter=1
WILKINSON, M. D. \textit{et al.} The FAIR Guiding Principles
for scientific data management and stewardship.
\textbf{Scientific Data}, v. 3, 2016.
DOI: 10.1038/sdata.2016.18.\par\smallskip
\end{singlespace}

\vspace{1cm}

\noindent\textbf{Citar como:}\\
CHALHOUB, L. Cirurgia uroginecológica no Sistema Único de Saúde,
2008\textendash 2025: ganhos silenciosos de eficiência,
desigualdade territorial persistente, choque pandêmico e represa
cirúrgica. \textit{Mirante dos Dados
\textemdash{} Working Paper n. 5}, v. 1.0, abr. 2026.
Disponível em:
\url{https://leonardochalhoub.github.io/mirante-dos-dados-br/}.

\vspace{0.5cm}

\noindent\textbf{Reconhecimento:} este Working Paper tem como
base o trabalho de especialização em Enfermagem de TATIELI
(2022). A presente versão amplia o recorte temporal, incorpora
a correção da camada \textit{silver} (\textit{commit}
\texttt{fa869cf}) e introduz a leitura longitudinal de eficiência
e represa cirúrgica.

\vspace{0.5cm}

\noindent\textbf{Licença:} dados consolidados (camada \textit{gold})
e código do \textit{pipeline} distribuídos sob licença MIT.
Texto deste artigo distribuído sob licença Creative Commons
Atribuição 4.0 Internacional (CC BY 4.0).

\end{document}
